\section{Clean Customer Deployment and Complete Mother Migration}
\label{sec:pilot-customer-deployment-migration}

Pilot now turns the customer-ownership architecture into signed distributable and
migration artifacts.  This completes coded tasks \texttt{MOB-1001},
\texttt{MOB-1003} and \texttt{MOB-1007}.  A personal computer, household server,
private VPS or business deployment remains owned by the customer; Tessaris need not
receive a readable copy of its intelligence or data.

\subsection{Clean deployment packages}

The mother-ownership authority builds the same application source for four explicit
profiles: personal computer, home server, private VPS and business.  The generated
configuration selects a loopback or network binding appropriate to that profile and
places writable memory under \texttt{customer-data} and credentials under
\texttt{customer-secrets}.  Neither directory is included in the distributable.

Packaging rejects symbolic links, over-large sources and an output nested inside its
input.  It excludes runtime state, version-control state, build caches,
\texttt{node\_modules}, environment files and filenames that imply passwords, tokens,
API keys, cookies, credentials or private keys.  Every included application file has
an exact byte count and SHA-256 digest.  The release contains the application,
profile configuration, plain-language instructions and a launcher, plus an
Ed25519-signed canonical manifest.  Verification rejects an altered signature,
configuration or application file before installation is trusted.

\subsection{Connect to my Pilot}

The unified browser application already implements the customer-owned connection
flow.  A user enters the HTTPS address displayed by their mother.  The browser reads
the mother's well-known descriptor, verifies its signature and endpoint, compares
human-readable trust words, and completes a six-digit local confirmation.  The phone
creates and retains its own signing key and receives a persona-bound certificate and
short-lived lease.  The flow connects to the selected customer mother directly; it
does not authenticate through, or upload mother secrets to, a Tessaris account.

After pairing, requests prefer the original local address, then a mother-signed
customer remote endpoint, and finally an optional opaque encrypted relay.  Every
candidate must prove the original mother identifier and key fingerprint.  A network
failure may change transport; an application or policy denial may not.  This closes
the coded browser-access task, while public DNS, certificates and firewall policy
remain deployment-specific operator work.

\subsection{Encrypted complete-mother migration}

A migration export archives the operator-selected complete mother directory, records
the path, byte count and digest of every file, then encrypts the archive with
AES-256-GCM.  Its key is derived from a user passphrase using scrypt.  The encrypted
backup is placed inside a two-file migration capsule with a signed manifest binding
the source mother identifier and fingerprint, export identifier, creation time,
encrypted payload digest and file count.

Inspection verifies the source-mother signature and encrypted payload digest without
decrypting customer content.  Restore verifies the passphrase and authenticated
ciphertext, rejects unsafe archive paths, extracts the files, regenerates the
manifest and compares it with the source record.  The receipt explicitly records
that identity, permissions, messages and proof state were preserved.

Migration does not silently make the destination writable.  The capsule records that
the source primary must be released and the restored destination must acquire the
existing signed single-primary lease.  This makes the cut-over visible and prevents
a successful file copy from being confused with authority to run two writable
mothers.  A physical two-host rehearsal and rollback exercise remains a production
closure gate.

\subsection{Customer-self-hosted opaque relay}

Advanced users and businesses can run the reference relay as a separate TLS 1.2 or
newer service in their own VPS or private cloud.  The service exposes health,
signed-route registration, encrypted request submission and polling, encrypted
response submission and retrieval, and a minimum operator status surface.  A route
is registered only after verifying the issuing mother's descriptor signature.

The relay holds no payload-decryption key.  Phone submission and mother polling use
different high-entropy capabilities, stored by the queue as canonical hashes.  The
service validates envelope structure, timestamp, expiry, nonce and ciphertext size;
deduplicates request identifiers; bounds pending work per route; and applies a
per-address request limit.  HTTPS responses disable caching, framing and content
sniffing and enable HSTS.  The command-line entry point requires an operator-supplied
certificate chain and private key, so Tessaris is not placed into the customer's
network or certificate authority.

The queue is deliberately ephemeral because relay envelopes live for at most two
minutes and the mother remains authoritative.  A relay restart can require route
re-registration and request retry under the existing idempotency contract; it cannot
destroy customer memory or messages held by the mother.  A production operator may
place multiple instances behind its own availability layer, but must not add readable
logging of bodies, tokens or decrypted content.

\subsection{Verification and remaining Stage 10 work}

Automated tests build and verify all four deployment profiles, prove that private
runtime and credential fixtures do not ship, restore identity, permission, message
and proof fixtures byte-for-byte, reject a wrong migration passphrase, and retain the
existing split-brain tests.  Live HTTPS tests exercise the customer-self-hosted relay
through a complete encrypted request/response round trip and reject unsigned route
registration.  The browser connection tests verify direct mother
discovery, signed identity pinning, local confirmation, customer-owned direct-route
failover and opaque-relay fallback.

All coded Stage 10 tasks are now implemented.  Production installation, public
certificate, firewall, relay-restart and two-host migration exercises remain field
gates rather than being inferred from unit tests.
